% Versi cetak contoh/PANDUAN_METODE.md (gaya laporan praktikum). Ubah md dulu, lalu samakan di sini.
% Bangun: pdflatex PANDUAN_METODE.tex (dua kali), hapus .aux/.log/.out.
\documentclass[14pt,a4paper]{extarticle}
\usepackage[left=3cm,right=2.5cm,top=2.5cm,bottom=2.5cm]{geometry}
\usepackage[T1]{fontenc}
\usepackage[utf8]{inputenc}
\usepackage[indonesian]{babel}
\usepackage{lmodern}
\usepackage{booktabs,tabularx,xltabular,array,enumitem,titlesec,caption,float}
\usepackage{url}
\usepackage[hidelinks]{hyperref}
\setlength{\parindent}{1cm}
\setlength{\parskip}{2pt}
\setlength{\emergencystretch}{2em}
\titleformat{\section}{\normalfont\bfseries}{\thesection.}{0.5em}{}
\titleformat{\subsection}{\normalfont\bfseries}{\thesubsection}{0.5em}{}
\titlespacing*{\section}{0pt}{8pt}{2pt}
\titlespacing*{\subsection}{0pt}{6pt}{1pt}
\setlist{leftmargin=*,itemsep=0pt,parsep=0pt,topsep=1pt}
\captionsetup{font=small,labelfont=bf,skip=3pt}
\setlength{\textfloatsep}{6pt}
\setlength{\intextsep}{6pt}
\setlength{\LTpre}{4pt}
\setlength{\LTpost}{4pt}
\newcolumntype{Y}{>{\raggedright\arraybackslash}X}
\newcommand{\code}[1]{\texttt{\detokenize{#1}}}
\urlstyle{tt}
\makeatletter\def\UrlBreaks{\do\_\do/\do.}\makeatother
\newcommand{\nb}[1]{\path{#1}}
\newcommand{\lbl}[1]{\textbf{#1}}

\begin{document}

\begin{center}
  {\bfseries PANDUAN METODE NICHE ARSENAL TEKS}\\
  {\bfseries GEMASTIK XIX Penambangan Data}\\
  {Kapan dipakai dan untuk data seperti apa}
\end{center}

Tujuh notebook di \nb{contoh/text/} pada panduan ini adalah lapis kedua. Notebook ini tidak
menggantikan baseline dan baru dipakai setelah model dasar berjalan dan menulis OOF
(\nb{klasifikasi_tfidf}, \nb{baseline_oof}, \nb{teks_tabular}, \nb{finetune_oof}). Peta soal ke
notebook ada di \nb{contoh/PANDUAN.md}.

\section{Aturan Main}
Aturan berikut berlaku untuk ketujuh notebook.
\begin{enumerate}
  \item \lbl{Pakai hanya kalau skor jujur naik.} Tiap notebook mencetak skor jujur (lipatan luar
  atau holdout) dan memilih cara paling sederhana yang selisihnya $\leq$ \nb{MIN_GAIN} dari yang
  terbaik. Ikuti pilihan itu. Di demo, cara rumit sering menang di OOF lalu kalah di test.
  \item \lbl{Semuanya untuk target kelas.} Target angka memakai \nb{ensemble_oof},
  \nb{stacking_teks_tabular_regresi}, dan \nb{interval_prediksi}.
  \item \lbl{Metrik harus ``makin besar makin baik''.} Untuk logloss, RMSE, dan MAE, tulis
  \code{return -(...)}.
  \item \lbl{Hapus sel bertanda \code{HAPUS saat lomba}.} Yang paling berbahaya adalah sel ``OOF
  demo'' di \nb{bobot_ensemble_metaheuristik} dan \nb{stacking_fwls} yang berisi
  \code{OOF_FOLDER = OUTPUT_FOLDER}. Kalau lupa dihapus, notebook melatih model demo dan membaca
  OOF itu, bukan OOF milik sendiri.
  \item Data demo kecil (emot 3,9 ribu baris, baju 17,6 ribu baris). Selisih $<$ 0,005 pada data
  sekecil itu masih bisa berupa derau. Kemenangan penting dicek dengan \nb{uji_beda_model}.
  \item Batas angka seperti ``$\geq\pm$5.000 baris'' atau ``kelas $\leq\pm$10'' adalah perkiraan,
  bukan hasil uji. Skor jujur di notebook tetap yang menentukan.
\end{enumerate}

\section{Pilih Cepat}
Tabel~\ref{tab:pilih} memetakan keadaan data atau model ke notebook yang dibuka.
{\small
\begin{xltabular}{\linewidth}{@{}Y>{\raggedright\arraybackslash}p{4.6cm}>{\raggedright\arraybackslash}p{3.4cm}@{}}
\caption{Keadaan data dan notebook yang dibuka}\label{tab:pilih}\\
\toprule
Keadaan data / model & Buka & Syarat minimum\\
\midrule
\endfirsthead
\toprule
Keadaan data / model & Buka & Syarat minimum\\
\midrule
\endhead
\bottomrule
\endlastfoot
$\geq$ 3 OOF, metrik berundak (F1 makro, akurasi, QWK), ingin bobot lebih baik dari rata-rata
 & \nb{bobot_ensemble_metaheuristik} & file \code{oof_}/\code{test_} dari lipatan yang sama\\
 \addlinespace[3pt]
Kualitas anggota timpang, atau model unggul di jenis baris berbeda (teks pendek vs panjang, kolom
kosong) & \nb{stacking_fwls} & $\geq$ 2 OOF dan kolom angka untuk ciri baris\\
 \addlinespace[3pt]
Ensemble mentok karena semua anggota TF-IDF + linear (korelasi $>$ 0,95)
 & \nb{pengklasifikasi_alternatif} & teks\\
 \addlinespace[3pt]
Teks + tabel untuk LightGBM, tidak ada waktu merancang fitur & \nb{fitur_otomatis_tetangga}
 & kolom angka/kategori\\
 \addlinespace[3pt]
Fitur buatan banyak ($>$ 30), curiga derau, perlu slide ``fitur penting''
 & \nb{seleksi_fitur_metaheuristik} & target biner\\
 \addlinespace[3pt]
Label sedikit, teks tanpa label banyak & \nb{co_training_teks} (juga \nb{graf_label_teks},
 \nb{semi_supervised_teks}) & teks $\geq$ beberapa kata\\
 \addlinespace[3pt]
Juri butuh penjelasan ``JIKA \dots\ MAKA'' & \nb{rulefit_aturan} & biner (multikelas: 1 lawan
 sisa)\\
\end{xltabular}}

\section{Kartu per Notebook}
Setiap kartu berisi data yang cocok, kapan notebook tidak dipakai, cara membaca hasil, hasil demo,
dan bagian yang perlu diganti saat lomba.

\subsection{Bobot Gabungan tanpa Turunan}
Notebook: \nb{bobot_ensemble_metaheuristik}.
\begin{itemize}
  \item \lbl{Bagus untuk:} klasifikasi dengan metrik berundak (F1 makro, akurasi, QWK) yang
  membuat Nelder-Mead/SLSQP macet; 4--8 model yang skornya berdekatan; train $\geq\pm$5.000 baris.
  \item \lbl{Jangan dipakai kalau:} hanya 2 model (cukup rata-rata di \nb{ensemble_oof}); satu
  model jauh lebih unggul (bobot lari ke satu model); metrik halus seperti logloss/AUC (cukup
  Nelder-Mead di \nb{ensemble_oof}); data $<$ 2.000 baris (bobot cocok ke derau).
  \item \lbl{Pilih cara:} lihat Tabel~\ref{tab:bobot}.
  \item \lbl{Baca:} kolom \code{skor jujur}, bukan \code{skor OOF (optimis)}. Selisih keduanya
  menunjukkan besar overfit bobot.
  \item \lbl{Demo (emot, 5 model, 1 menit):} rata-rata 0,675 naik menjadi 0,685 (skor jujur hill
  negatif). Di test justru hill biasa terbaik (0,674 vs 0,664). Untuk data $<$ 5.000 baris,
  naikkan \nb{MIN_GAIN} ke 0,005.
  \item \lbl{Ganti:} \nb{MODELS}, \nb{metric}, dan hapus sel OOF demo.
  \item \lbl{Untuk juri:} ``F1 makro tidak bisa diturunkan, jadi bobot dicari dengan
  metaheuristik dan dinilai di lipatan yang tidak dipakai mencari bobot.''
\end{itemize}
\begin{table}[H]\centering\small
\caption{Pilihan cara mencari bobot}\label{tab:bobot}
\begin{tabularx}{\linewidth}{@{}lY@{}}
\toprule
Cara & Pakai kalau\\
\midrule
Hill (Caruana) & bawaan; paling tahan overfit karena bobot berupa hitungan langkah\\
Hill negatif & ada model yang salahnya bisa ``dikurangkan''; butuh data besar\\
Annealing & hill macet di bobot yang sama\\
Evolusi diferensial / genetik & lebih dari 6 model, ruang bobot besar\\
Optuna & perlu jejak percobaan; hasilnya mirip evolusi\\
\bottomrule
\end{tabularx}
\end{table}

\subsection{Bobot yang Berubah Menurut Ciri Baris (FWLS)}
Notebook: \nb{stacking_fwls}.
\begin{itemize}
  \item \lbl{Bagus untuk:} teks + kolom angka; panjang teks sangat beragam; ada kolom teks yang
  sering kosong; anggota punya kekuatan berbeda (huruf, kata, judul saja, tabular). Train
  $\geq\pm$5.000 baris karena koefisien FWLS berjumlah model $\times$ kelas $\times$ ciri.
  \item \lbl{Jangan dipakai kalau:} anggota mirip semua; data kecil; tidak ada ciri baris yang
  bermakna.
  \item \lbl{Baca:} \nb{skor_per_panjang.png}. Garis model yang bersilangan antar kelompok panjang
  berarti FWLS layak. Garis yang sejajar berarti stacking LR biasa sudah cukup.
  \item \lbl{Demo (baju rating 5 kelas, 4 model, $\pm$1,5 menit):} hasil pada
  Tabel~\ref{tab:fwls}. Rata-rata merusak karena model tabular lemah (0,159). Kalau anggota
  timpang, stacking apa pun jauh di atas rata-rata. FWLS sendiri tidak menang di demo ini.
  \item \lbl{Ganti:} \nb{MODELS}, \nb{TEXT_COLS}, \nb{META_COLS} (kolom yang menggambarkan
  \textit{jenis} baris, bukan yang paling kuat menebak target), dan hapus sel OOF demo.
\end{itemize}
\begin{table}[H]\centering\small
\caption{Hasil demo \code{stacking_fwls} (F1 makro)}\label{tab:fwls}
\begin{tabularx}{\linewidth}{@{}Ycc@{}}
\toprule
Cara & Skor jujur & Skor test\\
\midrule
Model tunggal terbaik & 0,434 & 0,435\\
Rata-rata & 0,357 & 0,359\\
\textbf{Stacking LR (dipilih)} & \textbf{0,496} & \textbf{0,509}\\
FWLS & 0,493 & 0,490\\
LightGBM + ciri & 0,483 & 0,497\\
\bottomrule
\end{tabularx}
\end{table}

\subsection{Anggota Ensemble yang Salahnya Berbeda}
Notebook: \nb{pengklasifikasi_alternatif}.
\begin{itemize}
  \item \lbl{Bagus untuk:} teks saja, ketika ensemble sudah ada tetapi semua anggotanya TF-IDF +
  linear. Teks pendek penuh typo/alay cocok dengan gzip-kNN dan kernel RFF. Dokumen panjang yang
  mirip pencarian cocok dengan BM25-kNN dan Rocchio. Kelas tidak seimbang cocok dengan Complement
  NB.
  \item \lbl{Jangan dipakai kalau:} sudah ada fine-tune transformer yang kuat (tambahan jarang
  membantu); data $>$ 20 ribu baris untuk gzip (lambat).
  \item \lbl{Baca:} kolom \code{+ patokan} adalah skor rata-rata model itu dengan patokan; ambil
  yang di atas patokan. Kolom \code{korelasi} makin kecil berarti makin beragam. Kolom
  \code{menyelamatkan} adalah porsi baris yang salah di patokan tetapi benar di model ini.
  Keputusan akhir diambil di \nb{ensemble_oof} atau \nb{bobot_ensemble_metaheuristik}.
  \item \lbl{Demo (emot, 1,5 menit):} sendirian semuanya kalah (gzip 0,365, terbaik
  passive-aggressive 0,675 vs patokan 0,679). Dirata-rata dengan patokan, OOF naik (RFF 0,687,
  Rocchio 0,684) tetapi test tidak. Perlakukan sebagai calon anggota, bukan model utama.
  \item \lbl{Waktu:} gzip sebanding dengan baris $\times$ \nb{NCD_POOL} (3,9 ribu baris $\approx$ 1
  menit, 20 ribu $\approx$ 5 menit). Kecilkan \nb{NCD_POOL} kalau lambat.
  \item \lbl{Hasil:} \code{oof_alt_*.npy} di \nb{OOF_FOLDER}, langsung ditulis di \nb{MODELS}
  notebook ensemble.
\end{itemize}

\subsection{Fitur Tetangga kNN dan OpenFE Mini}
Notebook: \nb{fitur_otomatis_tetangga}.
\begin{itemize}
  \item \lbl{Bagus untuk:} teks + tabel untuk LightGBM, dengan kategori bernilai banyak dan kolom
  angka yang bisa digabung (rasio, selisih terhadap rata-rata kelompok).
  \item \lbl{Jangan dipakai kalau:} teks saja; sudah memakai \nb{stacking_teks_tabular} (sudah
  punya fitur tetangga, jadi ambil kolom otomatisnya saja).
  \item \lbl{Baca:} \nb{banding_fitur.csv} (dasar, + tetangga, + otomatis, + keduanya). Untung
  fitur otomatis yang kecil itu normal.
  \item \lbl{Demo (baju biner, 30 detik):} AUC 0,906 naik menjadi 0,914 dengan fitur tetangga dan
  0,916 dengan keduanya. Fitur otomatis terbaik adalah \code{title_huruf/title_kata} (rata-rata
  panjang kata judul).
  \item \lbl{Pakai hasilnya:} \nb{fitur_baru_train.csv} dan \nb{fitur_baru_test.csv} (ada
  \nb{ID_COL}, urutan baris sama dengan file asli), lalu \code{pd.concat} kolomnya ke notebook
  LightGBM.
  \item \lbl{Belum diuji:} multikelas (kodenya ada).
\end{itemize}

\subsection{Membuang Fitur Derau}
Notebook: \nb{seleksi_fitur_metaheuristik}.
\begin{itemize}
  \item \lbl{Bagus untuk:} lebih dari 30 fitur buatan (SVD, statistik teks, encoding, agregat);
  $\geq$ 5.000 baris; ingin model lebih ramping atau cepat; bahan slide ``fitur mana yang penting
  dan bukti sisanya derau''.
  \item \lbl{Jangan dipakai kalau:} fitur $<$ 20; ribuan kolom TF-IDF (terlalu lambat, pakai
  chi\textsuperscript{2} atau L1); target multikelas atau angka (notebook ini biner saja); sisa
  waktu $<$ 30 menit.
  \item \lbl{Pilih cara:} lihat Tabel~\ref{tab:seleksi}.
  \item \lbl{Demo (baju biner, 48 fitur termasuk 10 kolom acak, 1 menit):} semua cara membuang 10
  kolom acak, tetapi holdout tidak naik (semua fitur 0,9066, permutasi 0,9067), sehingga notebook
  tetap memilih semua fitur. LightGBM sudah tahan derau, jadi gunanya terutama untuk penjelasan
  dan kecepatan.
\end{itemize}
\begin{table}[H]\centering\small
\caption{Pilihan cara seleksi fitur}\label{tab:seleksi}
\begin{tabularx}{\linewidth}{@{}lY@{}}
\toprule
Cara & Sifat\\
\midrule
Null importance & cepat; grafiknya bagus untuk slide\\
Boruta & paling ketat (membuang paling banyak)\\
Permutasi & paling lunak\\
Annealing / genetik & mencari kombinasi; paling mahal, untuk 30--100 fitur\\
\bottomrule
\end{tabularx}
\end{table}

\subsection{Semi-supervised Dua Sudut Pandang}
Notebook: \nb{co_training_teks}.
\begin{itemize}
  \item \lbl{Bagus untuk:} label ratusan sampai $\pm$2 ribu, teks tanpa label jauh lebih banyak
  (test besar atau file tambahan), kelas $\leq\pm$10.
  \item \lbl{Jangan dipakai kalau:} label sudah banyak ($>\pm$50\% data) karena untungnya hilang;
  teks hanya 1--3 kata sehingga sudut pandang kata dan huruf tidak lagi saling melengkapi; test
  bergeser dari train (cek \nb{geser_distribusi}) karena label palsu yang salah akan menumpuk.
  \item \lbl{Pilih notebook:} lihat Tabel~\ref{tab:semi}.
  \item \lbl{Demo (emot, 15\% label = 544 baris, 25 detik):} hasil pada Tabel~\ref{tab:cotrain}.
  Ambang keyakinan sederhana mudah merusak; jumlah tetap per kelas lebih aman.
  \item \lbl{Ganti:} \code{LABEL_SHARE = 1.0} saat lomba (0,15 hanya untuk demo). Baris train
  bertarget kosong otomatis masuk kolam tak berlabel. File teks tambahan tanpa label diisi di
  \nb{UNLABELED_CSV}.
\end{itemize}
\begin{table}[H]\centering\small
\caption{Pilihan notebook semi-supervised}\label{tab:semi}
\begin{tabularx}{\linewidth}{@{}lY@{}}
\toprule
Notebook & Cocok kalau\\
\midrule
\nb{co_training_teks} & ada dua sudut pandang (kata vs huruf, teks vs tabel)\\
\nb{graf_label_teks} & teks membentuk kelompok/topik yang jelas, test besar\\
\nb{semi_supervised_teks} & kelas banyak; pseudo-label N teratas per kelas\\
\bottomrule
\end{tabularx}
\end{table}
\begin{table}[H]\centering\small
\caption{Hasil demo \code{co_training_teks} (F1 makro)}\label{tab:cotrain}
\begin{tabularx}{\linewidth}{@{}Ycc@{}}
\toprule
Cara & OOF berlabel & Skor test\\
\midrule
Hanya label & 0,548 & 0,557\\
Self-training & 0,493 & 0,511\\
\textbf{Co-training (dipilih)} & \textbf{0,569} & \textbf{0,578}\\
Tri-training & 0,519 & 0,540\\
\bottomrule
\end{tabularx}
\end{table}

\subsection{Aturan ``JIKA \dots\ DAN \dots\ MAKA'' untuk Slide}
Notebook: \nb{rulefit_aturan}.
\begin{itemize}
  \item \lbl{Bagus untuk:} presentasi ke juri non-teknis; target biner atau satu kelas penting;
  campuran kata kunci dan kolom tabel.
  \item \lbl{Jangan dipakai kalau:} sebagai model utama (CV di bawah LightGBM); multikelas dengan
  banyak kelas (buat satu per kelas penting saja).
  \item \lbl{Baca:} \nb{aturan_teratas.png}. Biru menaikkan peluang, oranye menurunkan. Aturan
  ``TIDAK kata:x DAN TIDAK kata:y'' berarti kata-kata negatif tidak muncul, biasanya ulasan
  positif. Untuk slide, pilih 3--5 aturan dengan \code{porsi baris} 5--60\%. Aturan yang berlaku
  untuk 95\% baris kurang bercerita.
  \item \lbl{Demo (baju biner, 22 detik):} AUC CV RuleFit 0,873, LightGBM 0,877, L1 0,879; 53
  aturan terpakai dari 556; test 0,881.
  \item \lbl{Ganti:} \nb{POSITIVE}, \nb{STOP_WORDS} (\code{None} atau daftar stopword
  Indonesia), dan \nb{RULE_C} (makin kecil, aturan makin sedikit).
\end{itemize}

\section{Alur Gabungan yang Disarankan}
\lbl{Teks saja, multikelas.}
\begin{enumerate}
  \item TF-IDF dan fine-tune menulis OOF.
  \item \nb{pengklasifikasi_alternatif} untuk mencari anggota yang berbeda.
  \item \nb{ensemble_oof}.
  \item Kalau metrik F1/QWK dan ada $\geq$ 4 model: \nb{bobot_ensemble_metaheuristik}.
  \item \nb{uji_beda_model} untuk memastikan kemenangan bukan derau.
\end{enumerate}
\lbl{Teks + tabel, biner.}
\begin{enumerate}
  \item \nb{teks_tabular} atau \nb{stacking_teks_tabular}.
  \item \nb{fitur_otomatis_tetangga}.
  \item \nb{seleksi_fitur_metaheuristik}.
  \item LightGBM final.
  \item \nb{rulefit_aturan} dan \nb{penjelasan_teks} untuk slide.
\end{enumerate}
\lbl{Teks + tabel, multikelas, anggota timpang.} Beberapa OOF digabung di \nb{stacking_fwls}.

\noindent\lbl{Label sedikit.} \nb{co_training_teks} dan \nb{graf_label_teks}, lalu OOF keduanya
digabung di \nb{ensemble_oof}.

\section{Urutan saat Waktu Sempit}
Tabel~\ref{tab:urutan} mengurutkan langkah menurut untung di demo.
\begin{table}[H]\centering\small
\caption{Urutan langkah menurut untung demo}\label{tab:urutan}
\begin{tabularx}{\linewidth}{@{}cY>{\raggedright\arraybackslash}p{2.6cm}>{\raggedright\arraybackslash}p{1.9cm}>{\raggedright\arraybackslash}p{2.6cm}@{}}
\toprule
No & Langkah & Untung demo & Waktu demo & Kapan\\
\midrule
1 & Stacking LR saat anggota timpang (\nb{stacking_fwls} atau \nb{ensemble_oof}) & +0,06 F1 vs
 model terbaik & $\pm$1,5 menit & selalu dicoba\\
2 & Fitur tetangga (\nb{fitur_otomatis_tetangga}) & +0,008 AUC & 30 detik & selalu dicoba\\
3 & Co-training & +0,02 F1 & 25 detik & hanya kalau label sedikit\\
4 & Bobot metaheuristik, seleksi fitur, pengklasifikasi alternatif & kecil/tidak stabil & 1--1,5
 menit & sisa waktu $>$ 1 jam\\
5 & RuleFit & -- & 22 detik & khusus bahan slide\\
\bottomrule
\end{tabularx}
\end{table}
Semua waktu di atas diukur di laptop dengan data demo (4--18 ribu baris). Data 5$\times$ lebih
besar kira-kira 5$\times$ lebih lama; gzip-kNN serta annealing/genetik bisa lebih lama lagi.

\end{document}
